% Source: https://artofproblemsolving.com/community/c3h1564870_gcdmods
% Stable identity: aops-topic-1564870-q01
\textbf{The statement is false.} Take $a=9\equiv1\pmod4$ and $b=3\equiv3\pmod4$: $\gcd(9,3)=3$. (Another example: $a=5$, $b=15$, with $\gcd=5$.) In fact for any odd $d$, the numbers $a=d^2\equiv1\pmod 4$ and $b=d$ or $b=3d$ (whichever is $\equiv3\pmod4$) share the factor $d$.

\textbf{What is true.} $\gcd(a,b)$ is always \emph{odd}: $a$ and $b$ are both odd, so $2$ divides neither, and every common divisor is odd. Moreover every odd number $d$ occurs as such a gcd, by the construction above (for $d=1$, take $a=1$, $b=3$). So nothing more than ``$\gcd(a,b)$ is odd'' can be concluded. $\blacksquare$
